\heading{数据结构与算法A-15秋-期末}

\noteinfo{本卷由原试卷 OCR 精修；题面、参考答案和必要推导均整理为可维护的 TeX。}

\textbf{一、选择题（每题 2 分，共 20 分）}

\begin{question}
若要一个运行时间为 $100n^2$ 的算法在相同机器上快于另一个运行时间为 $2n$ 的算法，则最小的 $n$ 值应为 \blank。
\begin{solution}
若按合理题意理解为比较 $100n^2$ 与 $2^n$，最小 $n$ 为 $15$。检验临界值：$n=14$ 时 $100n^2=19600$，而 $2^{14}=16384$，指数项还没有超过；$n=15$ 时 $100n^2=22500$，而 $2^{15}=32768$，首次超过。若严格按题面中的 $2n$ 理解，则线性函数 $2n$ 不可能在正整数范围内超过 $100n^2$，因此题目应有排版或 OCR 误差。
\end{solution}
\end{question}

\begin{question}
某算法仅含顺序执行的语句 1 和语句 2，语句 1 执行次数为 $20n^2+3n\log n$，而语句 2 执行次数为 $2n^3$，则该算法的时间复杂度为 \blank。
\begin{solution}
时间复杂度为 $O(n^3)$。顺序执行的两段代码总执行次数相加为 $20n^2+3n\log n+2n^3$。当 $n$ 足够大时，三次项增长最快，低阶项和常数系数在大 $O$ 记号中忽略，因此总时间复杂度为 $O(n^3)$。
\end{solution}
\end{question}

\begin{question}
下面术语中，\blank 与数据结构的存储结构无关。
\begin{choiceoptions}[5]
  \item 顺序表
  \item 链表
  \item 队列
  \item 循环链表
  \item 堆
\end{choiceoptions}
\begin{solution}
与存储结构无关的是 C、E。顺序表、链表、循环链表都直接描述元素在存储器中的组织方式；队列描述队尾入队、队首出队的操作约束，属于逻辑结构/抽象数据类型。堆强调完全二叉树上的偏序关系，通常可用数组实现，但概念本身不是某一种存储结构。
\end{solution}
\end{question}

\begin{question}
在一个具有 $n$ 个结点的有序单链表中插入一个新结点并仍保持其有序的时间复杂度为 \blank。
\begin{choiceoptions}[2]
  \item $O(n\log_2 n)$
  \item $O(1)$
  \item $O(n)$
  \item $O(n^2)$
\end{choiceoptions}
\begin{solution}
选 C，时间复杂度为 $O(n)$。有序单链表在已知前驱结点时插入只需 $O(1)$ 修改指针，但为了保持有序，必须从表头顺链比较关键字直到找到插入位置。平均检查约一半结点，最坏检查全部结点，因此时间阶为 $O(n)$。
\end{solution}
\end{question}

\begin{question}
若元素 $a$、$b$、$c$、$d$、$e$、$f$ 依次进栈，允许进栈、退栈操作交替进行，但不允许连续三次进行退栈操作，则不可能得到的出栈序列是 \blank。
\begin{choiceoptions}[2]
  \item dcebfa
  \item cbdaef
  \item bcaefd
  \item afedcb
\end{choiceoptions}
\begin{solution}
选 D。若第一个出栈元素为 $a$，必须先压入 $a$ 后立即弹出。随后要输出 $f,e,d,c,b$，需要继续压入 $b,c,d,e,f$，再按栈的逆序连续弹出 $f,e,d,\ldots$。其中 $f,e,d$ 已经构成连续三次出栈，违反题设限制，所以不可能。
\end{solution}
\end{question}

\begin{question}
表达式 $3 * 2 \char`\^ (4+2*2-6*3)-5$ 的求值过程中，当扫描到 $6$ 时，操作数栈和运算符栈为 \blank，其中 \char`\^ 为乘幂。
\begin{choiceoptions}[2]
  \item $3,2,4,1,1$；$(*\char`\^(+*-$
  \item $3,2,8$；$(*\char`\^-$
  \item $3,2,4,2,2$；$(*\char`\^(-$
  \item $3,2,8$；$(*\char`\^(-$
\end{choiceoptions}
\begin{solution}
选 D。括号内扫描到 $4+2*2$ 时，乘法先算出 $2*2=4$，再与前面的 $4$ 合成 $8$；遇到减号后，括号内当前等待的运算符为 $-$。外层的 $*$、乘幂和左括号尚未出栈，所以操作数栈为 $3,2,8$，运算符栈为 $(*\char`\^(-$。
\end{solution}
\end{question}

\begin{question}
有 4 棵结点存储整数关键码的二叉树，若它们的中序遍历序列分别为：A.~5, 3, 1, 2, 4；B.~1, 2, 3, 4, 5；C.~5, 4, 3, 2, 1；D.~1, 3, 5, 4, 2。请问其中可能是二叉搜索树的包括 \blank。
\begin{solution}
按通常二叉搜索树定义，左子树关键码小于根、右子树关键码大于根，因此中序遍历应得到关键码的非递减序列。四个候选序列中只有 $1,2,3,4,5$ 严格递增，因此选 B；其他序列都存在逆序对，例如先出现较大关键码再出现较小关键码，这会违反"中序遍历先访问左子树、再访问根、再访问右子树"的有序性，所以不可能是二叉搜索树的中序遍历。
\end{solution}
\end{question}

\begin{question}
（4 分）一棵 Huffman 树的带权外部路径长度定义为所有叶结点权值与其深度的乘积之和。设根结点的深度为 0，对集合 $\{ 2, 3, 5, 7, 8 \}$ 构造二叉 Huffman 树，则其最小带权外部路径长度为 \blank。
\begin{solution}
Huffman 合并时每次取当前最小的两个权值。先合并 $2$ 和 $3$ 得到 $5$，再合并两个 $5$ 得到 $10$，再合并 $7$ 和 $8$ 得到 $15$，最后合并 $10$ 和 $15$ 得到根 $25$。Huffman 树的 WPL 等于每次合并产生的新权值之和，因此
\[
WPL=5+10+15+25=55.
\]
这比逐个画出每个叶子的编码长度更快捷，但本质上等价于把每个叶权乘以其路径长度后求和。
\end{solution}
\end{question}

\begin{question}
两个字符串相等的充分必要条件是 \blank。（两字符串的长度相等且两字符串中对应位置的字符也相等）
\begin{solution}
两个字符串相等当且仅当长度相等且对应位置字符均相等。比较时通常先比较长度；若长度不同，即使较短串是较长串的完整前缀，也不是同一字符串。长度相同时，再从第一个字符开始逐位比较，只要存在某个位置字符不同，就可立即判定两个串不相等；只有所有位置都相同，才判定为相等。
\end{solution}
\end{question}

\begin{question}
2-3 树是一种满足以下两个条件的特殊的树：（1）每个内部结点有两个或三个子结点；（2）所有叶结点到根的路径长度相同。如果一棵 2-3 树有 9 个叶结点，那么它可能有 \blank 个非叶结点。
\begin{solution}
可能有 $4$ 或 $7$ 个非叶结点，取决于题目采用的高度和叶层口径。若根的三个孩子都是内部结点，并且这三个内部结点各有三个叶子，则叶结点数为 $3\times3=9$，非叶结点为根加三个孩子，共 $4$ 个。若允许在同一层继续把部分孩子细分为内部结点，同时仍保持所有叶子在同一层，则也可构造出 $7$ 个非叶结点。作答时应先说明采用的 2-3 树定义：每个内部结点有 2 或 3 个孩子，且所有叶子在同一层；再根据该定义数内部结点，不能只凭叶子数直接唯一确定非叶结点数。
\end{solution}
\end{question}

\textbf{二、辨析与简答题（共 36 分）}

\begin{question}
（6 分）什么是循环队列？循环队列的优点是什么？如何判别它的空和满？
\begin{solution}
循环队列把一维数组看成首尾相接，用取模运算让队尾到达数组末端后回到下标 $0$。它能克服普通顺序队列中"数组后端满但前端已有空位"的假溢出现象。若采用少用一个存储单元的约定，队空为 $front=rear$；入队前若 $(rear+1)\bmod maxSize=front$，则队满。也可额外维护元素个数，此时队空为 $count=0$，队满为 $count=maxSize$。
\end{solution}
\end{question}

\begin{question}
（6 分）在包含 $n$ 个结点的单链表、双链表和单循环链表中，若仅知道指针 $p$ 指向某结点，不知道表的头指针，能否将结点 $p$ 从相应的链表中删去？若可以，其时间复杂度各为多少？
\begin{solution}
单链表仅知 $p$ 通常不能真正删除 $p$，因为要修改前驱结点的 \texttt{next} 指针，而前驱无法由 $p$ 反向得到；若 $p$ 不是尾结点，可用"把后继结点的数据复制到 $p$，再删除后继结点"的折中办法，但这改变的是结点内容而非原结点身份。双链表可由 $p$ 的前驱/后继指针直接改链，时间 $O(1)$。单循环链表虽然能从 $p$ 出发绕一圈找到前驱，但最坏要扫描整表，时间 $O(n)$。
\end{solution}
\end{question}

\begin{question}
（8 分）请分别计算模式 $p = \text{``aabaac''}$ 优化前和优化后的 next 数组，并按照优化后的 next 数组针对目标 $t = \text{``aabaabaabaac''}$ 进行 KMP 快速模式匹配，请画出匹配过程的示意图并计算匹配过程中的比较次数。
\begin{solution}
优化后 next 向量为 $[-1,-1,1,-1,-1,2]$。求解时先按普通 KMP 的最长相等真前后缀得到回退位置，再检查若 $P[i]$ 与回退后的字符相同，则继续沿 next 回退，从而避免匹配时立即重复比较同一字符。用该向量匹配 $t=\texttt{aabaabaabaac}$ 时，前两段 \texttt{aabaab} 基本顺利推进，遇到末尾不匹配时按 next 跳转而不是回到主串下一位重新开始；逐次记录主串字符比较，可得比较次数为 $14$。
\end{solution}
\end{question}

\begin{question}
（8 分）根据树的双标记层次遍历序列，求其带度数的后根遍历序列。譬如，已知一棵树的双标记层次遍历序列如下：
\begin{ttbox}
A : ltag: 0 , rtag 1
B : ltag: 0 , rtag 0
C : ltag: 0 , rtag 1
D : ltag: 1 , rtag 0
G : ltag: 0 , rtag 1
E : ltag: 0 , rtag 1
H : ltag: 1 , rtag 1
F : ltag: 1 , rtag 0
I : ltag: 1 , rtag 1
\end{ttbox}
则其带度数的后根遍历序列为：D0 H0 G1 B2 F0 I0 E2 C1 A2（注：各个结点按照"结点名 度数"的方式给出，结点之间用空格分隔）。现某棵树的双标记层次遍历序列如下：
\begin{ttbox}
A : ltag: 0 , rtag 1
B : ltag: 0 , rtag 0
G : ltag: 1 , rtag 1
C : ltag: 0 , rtag 0
F : ltag: 0 , rtag 1
D : ltag: 1 , rtag 0
E : ltag: 1 , rtag 1
H : ltag: 1 , rtag 0
I : ltag: 1 , rtag 1
\end{ttbox}
请给出其带度数的后根遍历序列。
\begin{solution}
带度数后根遍历序列为 $D0\ E0\ C2\ H0\ I0\ F2\ B2\ G0\ A2$。求解时先由双标记层次序恢复父子关系：$ltag=0$ 表示有左子，$rtag=0$ 表示有右兄弟；然后对树做后根遍历，即先依次访问所有孩子子树，最后访问当前结点，并在结点名后写出孩子个数。
\end{solution}
\end{question}

\begin{question}
（8 分）使用重量权衡合并规则（权重合并规则）与路径压缩，对下列从 0 到 15 之间的数的等价对进行归并。在初始情况下，集合中的每个元素分别在独立的等价类中。当两棵树规模同样大时，使结点数值较大的根结点作为值较小的根结点的子结点。
\[
(0,2)\ (1,2)\ (3,4)\ (3,1)\ (3,5)\ (9,11)\ (12,14)\ (3,9)\ (4,14)\ (6,7)\ (8,10)\ (8,7)\ (7,0)\ (10,15)\ (10,13)
\]
请填写下面表格的空白部分树的父指针表示法的数组表示。也就是所有等价对都被处理之后，所得父结点的下标值（没有父结点则填"--1"）。

父结点下标

结点值

0 1 2 3 4 5 6 7 8 9 10 11 12 13 14 15

结点的下标 0 1 2 3 4 5 6 7 8 9 10 11 12 13 14 15

\begin{solution}
父指针数组为 $[-1,0,0,0,0,0,0,6,0,0,0,9,0,0,12,0]$。合并时始终把规模小的树挂到规模大的树根下；规模相同则根值较大的挂到根值较小的树下。路径压缩只在查找根的过程中改变沿途结点父指针，因此最终多数结点直接指向代表元 $0$，但 $7,11,14$ 等结点因最后一次压缩路径不同，仍分别指向 $6,9,12$。
\end{solution}
\end{question}

\textbf{三、算法填空题（每题 12 分，共 24 分）}

\begin{question}
（每空 3 分，共 12 分）下面的算法利用一个栈将一给定的序列从小到大排序。长度为 $n$ 的序列由 $1, 2, \dots, n$ 这 $n$ 个连续的正整数组成。请补全下面的代码段，使其可以判断给定的序列是否可以利用一个栈进行排序，如果可以，输出相应的操作。

例如：序列 4 3 1 2，此序列可以排序，输出：push push push pop push pop pop pop。

\begin{pycode}
class Stack:
    def clear(self): ...
    def push(self, item): ...
    def pop(self): ...
    def top(self): ...
    def is_empty(self): ...
    def is_full(self): ...


s = Stack()                  # s 为栈
sequence = [...]             # 待排序序列
n = len(sequence)


def is_legal():              # 判断序列是否可以排序
    for i in range(n):
        for j in range(i + 1, n):
            for k in range(j + 1, n):
                if ________:          # 填空 1
                    return False
    return True


def transform():             # 输出转换操作
    cur, i = 1, 0
    while ________:           # 填空 2
        while s.is_empty() or ________:    # 填空 3
            ________          # 填空 4
            print("push", end=" ")
            if cur == s.top():
                break
        s.pop()
        print("pop", end=" ")
        cur += 1
\end{pycode}
\begin{solution}
填空依次为：\texttt{\detokenize{sequence[k] < sequence[i] < sequence[j]}}；\texttt{\detokenize{cur <= n}}；\texttt{\detokenize{cur < s.top()}}；\texttt{\detokenize{s.push(sequence[i]); i += 1}}。第一空是在三重下标中判断是否出现严格递增的三元组；后面几空对应栈模拟扫描过程：\texttt{cur} 表示当前希望弹出的目标值，只有当还未越界且栈顶正好满足输出条件时才弹出，否则继续把输入序列元素入栈。这样既保证每个元素只入栈一次，也能检测给定序列是否可由合法栈操作产生。
\end{solution}
\end{question}

\begin{question}
（每空 3 分，分析 3 分，共 12 分）以下算法用来判断两棵树是否相同，试补全下面的代码段，并分析算法的运行时间代价。
\begin{pycode}
def is_equal(root1, root2):
    while root1 is not None and root2 is not None:
        if root1.value() != root2.value():
            return ______          # 填空 1
        if ______________________: # 填空 2
            return False
        root1 = root1.right_sibling()
        root2 = root2.right_sibling()
    if ________________:           # 填空 3
        return True
    return False
\end{pycode}
\begin{solution}
填空依次为：\texttt{\detokenize{False}}；\texttt{\detokenize{not is_equal(root1.left_most_child(), root2.left_most_child())}}；\texttt{\detokenize{root1 is None and root2 is None}}。递归比较两棵树是否相等时，首先处理一空一非空的情形，应直接返回 \texttt{False}；若两个根都为空，则说明对应子树同时结束，应返回 \texttt{True}。非空时先比较根信息，再递归比较最左孩子子树和右兄弟链；任一部分不相等就返回 \texttt{False}。每个结点只被访问常数次，故时间复杂度为 $O(n)$，递归栈深度为树高。
\end{solution}
\end{question}

\textbf{四、分析证明题（共 20 分）}

\begin{question}
（5 分）请证明非空满 $K$ 叉树的叶结点数目为 $(K-1)*n+1$，其中 $n$ 为分支结点数目。
\begin{solution}
用归纳法证明。$n=0$ 时没有分支结点，树只含一个根结点且它也是叶结点，叶数为 $1=(K-1)\cdot0+1$。假设有 $n$ 个分支结点时叶数为 $(K-1)n+1$。再增加一个分支结点，等价于把原来的某个叶结点扩展成有 $K$ 个孩子的内部结点：原叶减少 $1$ 个，新叶增加 $K$ 个，叶数净增 $K-1$。因此叶数变为 $(K-1)n+1+(K-1)=(K-1)(n+1)+1$。由归纳法，命题对所有 $n$ 成立。
\end{solution}
\end{question}

\begin{question}
（15 分）有 $n$ 个数 $K_1, K_2, \dots, K_n$ 顺序存储在数组中，形成一棵 $n$ 个结点的完全二叉树。现在要将它按如下方式建成一个最小值堆：从左至右扫描整个数组，将第 $i$ 个元素 $K_i$ $(i = 1, 2, \dots, n)$ 与其父结点 $K_{\lfloor i/2 \rfloor}$ 比较，如果 $K_i \ge K_{\lfloor i/2 \rfloor}$，则不再进行操作；若 $K_i < K_{\lfloor i/2 \rfloor}$，将 2 个结点对换，再将 $K_{\lfloor i/2 \rfloor}$ 与 $K_{\lfloor i/4 \rfloor}$ 比较，以此类推，直到比较到根结点 $K_1$。

\begin{enumerate}
  \item[(1)] 按照这种建堆方式，最坏情况下建堆的时间复杂度是多少？
  \item[(2)] 教材上使用的"筛选法"建堆（Siftdown），最坏情况下建堆的时间复杂度是多少？
  \item[(3)] 上述 2 种建堆法有类似之处，请问它们的时间复杂度是否一样，若不同，则请说明导致它们时间复杂度不同的原因。
\end{enumerate}
\begin{solution}
本题从左到右逐个上滤建堆，最坏情况下第 $i$ 层结点上滤 $i$ 次，总复杂度 $O(n\log n)$；筛选法自底向上下滤，越低层结点移动距离越短，总复杂度 $O(n)$。二者不同的根本原因是大量下层结点在上滤法中代价高，而在筛选法中代价低。
\end{solution}
\end{question}
